\documentclass[11pt]{article}

\usepackage[margin=1in]{geometry}
\usepackage{amsmath,amssymb,amsthm,mathrsfs}
\usepackage{booktabs,microtype}
\usepackage{xcolor}
\usepackage[colorlinks=true,linkcolor=blue!55!black,
  citecolor=green!45!black,urlcolor=blue!55!black]{hyperref}
\usepackage[nameinlink,noabbrev]{cleveref}
\numberwithin{equation}{section}

\newtheorem{theorem}{Theorem}[section]
\newtheorem{lemma}[theorem]{Lemma}
\newtheorem{proposition}[theorem]{Proposition}
\newtheorem{corollary}[theorem]{Corollary}
\newtheorem{conjecture}[theorem]{Conjecture}
\newtheorem{problem}[theorem]{Problem}
\theoremstyle{definition}
\newtheorem{definition}[theorem]{Definition}
\theoremstyle{remark}
\newtheorem{remark}[theorem]{Remark}

\newcommand{\Damp}{\mathcal D_{\mathrm{amp}}}
\newcommand{\Ample}{\mathcal A}
\newcommand{\Obs}{\operatorname{Obs}_{\mathrm{pc}}}

\title{Forbidden pc-Minors for Corner-Peelable Ample Partial Cubes}
\author{Olivier Bousquet}
\date{}

\hypersetup{
  pdftitle={Forbidden pc-Minors for Corner-Peelable Ample Partial Cubes},
  pdfauthor={Olivier Bousquet}
}

\begin{document}
\maketitle

\begin{abstract}
We study the forbidden partial-cube minors of corner-peelable ample
concept classes.  The nonample members are exactly the cubes with one
antipodal pair removed.  For the ample members, we give exact recognition
criteria and partial structural constructions. Using CCMW's equivalence
between corner peelings and acyclic unique-sink orientations, we interpret
rooted factors as prescribed-sink obstructions and separator gluing as
compatibility of outgoing boundary directions. Every ample obstruction
has a unique forbidden graph three-core, and finite labelled square
attachment data reconstruct all extensions of that core.  Obstructions
with a cut vertex have exactly two rooted factors; separating edges
admit an exact decomposition into two or three pieces. A finite-certificate
construction realizes three pieces in a $1237$-concept obstruction whose
nonempty-coordinate reductions all peel. Product attachments
convert peeling down to a prescribed ample subfamily into ordinary peeling
and transfer minimal relative failures to forbidden pc-minors. Replacing one
coordinate by a block, with its forbidden signs repeated across the block,
preserves forbiddenness and constructs obstructions with arbitrarily many
maximal reduction descent supports.  Compressing repeated signed-clause
columns gives a canonical reduced forbidden class with positive integer
multiplicities on its coordinates.  Graph pruning and column compression
alternate canonically; their inverse stages recover every obstruction from
a forbidden base of minimum degree at least three with distinct columns.
Every nonempty corner-peelable ample family occurs as a coordinate face of such a
base, even with isometric dimension minus VC dimension fixed at nine.
A further
construction replicates resolving additions to a fixed forbidden seed
over a layer cube.  For cornerless seeds, its terminal members are
parameterized by partitions of the signed layer facets.  Explicit
terminal examples show that this construction is not exhaustive, and
another infinite family rules out generation solely by corner deletions
from cubes.  The least cardinality of an ample obstruction lies between
\(64\) and \(267\).  A complete intrinsic generative classification
remains open: the missing ingredients lie in the intrinsic construction
and admissibility of these reduced forbidden bases and their factors.
\end{abstract}

\medskip
\noindent\textbf{Keywords.}
Partial cube, pc-minor, ample class, corner peeling, unique-sink orientation,
representation map, Yang complex,
simplicial ball, shellability, Alexander duality, linear quotients.

\setcounter{tocdepth}{2}
\tableofcontents
\clearpage

\section{Introduction}
\label{sec:damp-introduction}

A family of binary words \emph{shatters} a set of coordinates when its
projection onto those coordinates contains every binary pattern.
The family is \emph{ample} when each shattered set of coordinates also
supports a full cube inside the family.  These classes, also called
lopsided systems, were introduced
by Lawrence and developed from combinatorial and metric viewpoints
\cite{Lawrence83,BandeltCDK06}.  Join two words when they differ in one
coordinate.  The resulting graph is a partial cube: its graph distances
are exactly the Hamming distances between the words.

A vertex is a \emph{corner} if it belongs to a unique maximal cube.
A \emph{corner peeling} successively deletes corners until the family
is empty.  Reversing such a deletion order gives an ordering whose
nonempty prefixes are ample
\cite[Proposition~4.2]{ChalopiChepoiMoran22}.  We write \(\Damp\)
for the corner-peelable ample families.  The term \emph{dismantlable}
is also used for these families; it refers here to corner deletion,
rather than domination of vertices in an arbitrary graph.  Hall's
maximum class of VC dimension three has no corner and shows that an
ample class need not belong to \(\Damp\)
\cite[Theorem~4.5]{ChalopiChepoiMoran22}.

\subsection{The classification problem}

Partial-cube minors, or \emph{pc-minors}, combine convex restrictions
with contractions of coordinate classes \cite{KnauerMarc20}.  For
\(X\subseteq Q_E=\{0,1\}^E\), let \(\rho_e^iX\) be the section
with coordinate \(e\) fixed to \(i\), and let \(\pi_eX\) be the
projection deleting coordinate \(e\).  Constant coordinates are removed
from sections.  The class \(\Damp\) is pc-minor closed, so it is
determined by the basis
\[
 \Obs(\Damp)=\{O:O\notin\Damp,\;
       \text{every proper pc-minor of }O\text{ lies in }\Damp\}.
\]
Our aim is to understand the structure of this basis and give explicit
constructions of its members.

The nonample part is completely classified:
Proposition~\ref{prop:nonample-layer} identifies it as
\(\{Q_n\setminus\{x,\bar x\}:n\ge3\}\), where \(\bar x\)
is the antipode of \(x\).  For the ample part, pc-minor closure gives
the exact recognition criterion
\begin{equation}
 \begin{split}
 O\in\Obs(\Damp)\quad\Longleftrightarrow\quad
 &O\notin\Damp,\\[-2pt]
 &\rho_e^0O,\rho_e^1O,\pi_eO\in\Damp
       \quad(e\in E),
 \end{split}
 \label{eq:damp-introduction-recognition}
\end{equation}
for an ample family embedded with every coordinate active, meaning that
it occurs on an edge.
Equivalent algebraic and topological forms use linear quotients and
Yang-complex shellability
(Theorem~\ref{thm:damp-descendant-free-parametrization} and
Corollary~\ref{cor:damp-descendant-free-yang-ball-parametrization}).
Signed minimal nonfaces give a Boolean formulation
(Theorem~\ref{thm:damp-fixed-shadow-signed-clutter}).

An intrinsic generative description asks for more: finite structural
parameters, conditions specifying which parameters are admissible, a
construction that always produces an obstruction, and recovery of every
obstruction from such data.  The parameters may have arbitrary finite
size, and recursive constructions are allowed.  A bounded catalogue is
unnecessary.  For example, a down-set of arbitrary size can serve as a
structural parameter, as in the
\href{daisy_pc_minors.pdf}{companion classification for daisy cubes}.
For the present problem, a construction that assumes an arbitrary
forbidden seed still leaves that seed to be classified.  Likewise,
replacing the tests in \eqref{eq:damp-introduction-recognition} by
equivalent shelling or ordering conditions gives a recognition theorem
without supplying the desired structural inputs.

\subsection{Acyclic USOs and the organization of the results}
\label{sec:damp-intro-uso}

CCMW identify corner-peelable ample classes with those admitting an acyclic
unique-sink orientation (USO) \cite[Proposition~6.12]{ChalopiChepoiMoran22}.
For an ample class a USO requires both a unique sink on each contained cube
and containment of the cube spanned by each vertex's outgoing directions.
Its out-map records those directions and is a representation map
\cite[Theorems~6.1 and~6.8]{ChalopiChepoiMoran22}.
Section~\ref{sec:damp-uso-boundaries} gives the precise definitions.

This organizes the structural results by the information they preserve.
Canonical graph and signed-column reductions act on the support itself.
Rooted factors and separator constructions instead prescribe boundary data
on acyclic USOs of positive pieces. Peeling to a root means realizing it
as a sink. Peeling to a positive subfamily means extending its acyclic USO
with all boundary edges pointing inward
(Proposition~\ref{prop:damp-relative-uso-extension}). At a separating cube,
a boundary vertex can have outgoing exclusive directions in at most one
piece, because two would require an absent mixed square
(Proposition~\ref{prop:damp-uso-cube-gluing}). This single condition explains
the separator bound and the five boundary conditions used at an edge.

Negative certificates must exclude every acyclic USO. Protected missing
cubes do so by forcing a directed cycle in every possible USO
(Proposition~\ref{prop:damp-protected-uso-cycle}). One cyclic orientation
alone does not certify nonpeelability, and the existence of acyclic USOs
on proper minors does not make arbitrary inherited maps acyclic.
For retained targets that are themselves nonpeelable, relative peeling
requires only acyclicity outside the target; its full scope is preserved
in the deletion statements below.

\subsection{Canonical reduction and explicit constructions}

\paragraph{A unique core and all its square attachments.}
The strongest reduction applying to every ample obstruction is
Theorem~\ref{thm:damp-three-core}.  Repeatedly delete vertices of
degree less than three.  Each deletion is a degree-two corner deletion
that preserves forbidden-minor status.  The result is the unique graph
three-core \(K\), the largest induced subgraph of minimum degree at
least three.  It is nonempty, ample, and forbidden, and has the same
active coordinates as the original family.

The inverse operation is explicit.  Given an ample family, choose a
missing word and two coordinates such that the other three vertices
of its square are present and the coordinate pair is not shattered.
Adjoining that word is an ample square completion.  Over a forbidden
base, it preserves forbidden-minor status.  More generally, it preserves
the possibility of peeling to every retained subset of the base
(Proposition~\ref{prop:damp-rank-two-extension-invariance}).
To see the operation at
the smallest scale, \(\{00,01,10\}\) acquires its missing word
\(11\) using the coordinate pair \(\{1,2\}\); the three present
words specify the entire addition.

Theorem~\ref{thm:damp-square-attachment-data} turns a sequence of
these operations into finite structural data: abstract new vertices,
a coordinate pair and three targets for each vertex, and an acyclic
dependency digraph.  Coordinate consistency and injectivity determine
the new words; distinct pairs absent from the core's shattered sets
ensure that every addition is valid.  The data are recovered uniquely
from the resulting family, up to renaming the abstract new vertices.
Thus construction and recovery of \emph{all attachments} are settled.
The unresolved input is precisely which ample three-cores are forbidden.
Such a core can itself have corners.

\paragraph{Rooted factors and gluing.}
Say that \(A\) \emph{peels to \(a\)} when it has a corner peeling
ending at \(a\).  A rooted pair \((A,a)\) has the factor property
needed here exactly when
\begin{enumerate}
 \item \(A\) is ample and corner-peelable;
 \item no acyclic USO of \(A\) has sink \(a\);
 \item every proper pc-minor obtained while retaining the root, or
       its image under contraction, admits an acyclic USO with that sink.
\end{enumerate}
These are the \emph{rooted factors} used below.
Theorem~\ref{thm:damp-cut-vertex-obstructions} says that gluing any
two such factors at their roots gives a forbidden pc-minor.  Every
ample obstruction with a cut vertex arises this way.  It has exactly
two blocks and a unique cut vertex, so the factors are recovered
uniquely up to interchange.  This is a complete decomposition of that
branch, with an explicitly unresolved rooted-factor input.

The same inputs yield two-connected examples.  Put rooted factors
\(Y,C\) on disjoint coordinate sets, with common root \(0\), choose
an edge \(\{0,v\}\) of \(Y\), and choose a nonempty down-set
\(L\subseteq C\).  Here down-set means that replacing any ones by
zeros keeps a word in \(L\).  The family
\begin{equation}
 X=Y\cup C\cup(v+L)
 \label{eq:damp-introduction-downset}
\end{equation}
is forbidden, and is two-connected exactly when \(|L|>1\)
(Theorem~\ref{thm:damp-rooted-downset-completion}).  All unused
coordinates are zero, and \(v+L\) is the translated copy of \(L\)
across the chosen edge.  The marked coordinate blocks and edge recover
\(Y,C,L\).  The down-set is a structural parameter; the rooted
factors still satisfy the three recognition conditions above.

There is also exact recovery at a separating edge, meaning an edge
whose two endpoints form a vertex cut.  Its union with each component
outside the edge gives a canonical convex piece.  There are two or
three pieces, and Theorem~\ref{thm:damp-edge-separator-obstructions}
gives necessary and sufficient gluing conditions.  These use five attainable out-map conditions: an inward boundary at the
edge, either endpoint as sink, and either endpoint having just the common
edge as its outgoing edge. Equivalently, these specify retention of the
edge, its possible final endpoints, and its possible first endpoint
deletions as leaves.  Proposition~\ref{prop:damp-edge-boundary-transfer} realizes the boundary
profile $(0,1,1,0,0)$ from rooted inputs: both endpoint sinks are attainable,
but neither leaf exit nor edge retention is possible. Thus sink feasibility
alone cannot replace the five conditions. Theorem~\ref{thm:damp-three-piece-boundary-equivalence}
reduces existence of a three-piece forbidden union exactly to existence
of a positive piece failing edge retention but admitting both an endpoint
sink and an endpoint leaf exit. No minor-minimality assumption on that
witness is needed; both boundary properties survive edge-retaining minors.
Proposition~\ref{prop:damp-repair-face-rooted-factor} first gives a uniform
construction from a rooted factor, one repairing addition, and a proper
face containing the acquired cube. If that face omits a root neighbour,
the output is another rooted factor, with all proper rooted minor tests
automatically satisfied. A $308$-concept example results. The operation
has a recoverable shore presentation but is not closed under direct iteration;
Remark~\ref{rem:damp-repair-face-recovery} distinguishes this limitation
from coverage by the earlier ample-addition construction.
Theorem~\ref{thm:damp-three-piece-witness} supplies the other outcome with
profile $(0,1,0,1,0)$ on $330$ concepts, certified by explicit orders.
Gluing it, its reversal, and the $581$-concept piece gives a forbidden
$1237$-concept class with three components outside the edge. Thus the bound
of three is attained. All its nonempty-coordinate reductions peel, so this
example also belongs to the seed side of reduction descent.
A complementary backward step is given by
Proposition~\ref{prop:damp-combined-rooted-descent}: an ordinarily positive
nonroot-corner deletion followed by rooted minor minimalization reaches a
smaller rooted factor. If a proper minor is needed, it preserves an endpoint
repair on a face of the deleted corner cube. Iteration ends at factors whose
root is the unique source in every acyclic USO; whether these can have other
geometric corners remains a separate question. Corollary~\ref{cor:damp-joint-rooted-ordinary-descent}
bypasses that question by allowing a rooted deletion to enter the ordinary
nonpeelable branch. Its terminal inputs are rooted factors with a unique
geometric corner and cornerless ample forbidden pc-minors. Negative-minor
selection and admissible inverse expansions still prevent an intrinsic
forward generation theorem.
Compatibility between two pieces matters:
even in a forbidden three-core, an elementary minor can peel through
their interaction although one changed piece still fails to peel to
the shared edge (Remark~\ref{ex:damp-coupled-edge-minors-1162}).

The separator branch is also detectable from shattered pairs.  Its
crossing graph has one vertex per active coordinate and joins the
shattered pairs.  Disconnection detects cut vertices; after ruling those
out, a cut vertex of the crossing graph detects separating edges in
that direction.  Contracting the direction recovers all such edges.
If deleting that crossing-graph vertex leaves two components, there is
one separating edge and a unique pair of convex factors
(Proposition~\ref{prop:damp-crossing-separator-recovery} and
Corollary~\ref{cor:damp-shadow-edge-factorization}).  Thus the
remaining branch has a two-connected crossing graph; the peeling
properties of the recovered factors still require intrinsic criteria.

In the separator branch, descent through nonpeelable reductions
also has a canonical endpoint: there is at most one descent coordinate,
and after that one step all further reductions peel
(Theorem~\ref{thm:damp-separator-strong-descent}).  A nonpeelable
carrier must retain every remaining coordinate, so its direction is
adjacent to all others in the crossing graph.  Failure of two-connectivity
permits at most one such direction.  Whether its carrier actually peels
is still a separate condition.  Square attachments leave every
nonpeelable reduction unchanged as an embedded family
(Proposition~\ref{prop:damp-square-attachments-fixed-carriers}).
Hence the same conclusion holds whenever the canonical three-core has
a separator, even if the full obstruction has none.

\paragraph{Replacing coordinates by blocks.}
There is also a construction that applies to every ample forbidden input,
without choosing resolving additions.  Replace a coordinate \(z\) by a
nonempty block \(E\).  Keep its two sections at the constant words
\(0_E,1_E\), and put their union in every nonconstant section.
Theorem~\ref{thm:damp-coordinate-block-replacement} proves that the result
is forbidden exactly when the original family is forbidden.  On signed
minimal nonfaces, the rule simply replaces \(z\) by all of \(E\),
repeating its forbidden sign.  The two constant sections recover the input.
For \(|E|\ge2\), the output is its own graph three-core and has a
two-connected crossing graph.  If the input has no nonpeelable proper
reduction, the maximal descent supports are exactly
\(E\setminus\{e\}\), \(e\in E\), and every endpoint is the original
family with one coordinate renamed
(Corollary~\ref{cor:damp-coordinate-block-descent}).  Thus descent supports
need not be unique even though these examples share one terminal
isomorphism type.

The operation also has canonical unmarked recovery.  Group coordinates
whose columns in the signed minimal-nonface presentation agree up to sign,
and compress each group to one coordinate.  Ampleness forces the section
pattern required by the construction, and the reduced family is again
forbidden.  Theorem~\ref{thm:damp-canonical-signed-column-reduction}
proves that this gives an exact parametrization by reduced forbidden inputs
with distinct columns up to sign and arbitrary positive coordinate
multiplicities.  The description of these reduced inputs remains open.

The two canonical reductions also fit together.  Repeated columns force
minimum degree at least three, so graph pruning and column compression
are never available simultaneously in a forbidden class.  Exhaust one
type, then the other, until neither applies.
Theorem~\ref{thm:damp-alternating-reductions} proves that the resulting
base and all intermediate phases are canonical.  There are at most
\(2d-5\) phases for VC dimension \(d\), and the difference between
isometric dimension and VC dimension is unchanged.  Conversely, alternate
nonempty square attachment data and nontrivial coordinate multiplicities
over any forbidden base with minimum degree at least three and distinct
columns.  These are exactly the recovered construction stages; no further
compatibility tests are required.  Intrinsic admissibility of the final
forbidden bases remains the unresolved input.

This input does not become finite at a fixed value of the dimension
difference.  Theorem~\ref{thm:damp-reduced-base-universality} constructs,
from any nonempty ample family \(K\), a nonpeelable family \(X_K\)
with minimum degree at least four, distinct columns, complete crossing
graph, and dimension difference nine, such that
\[
 X_K\in\Obs(\Damp)\quad\Longleftrightarrow\quad K\in\Damp.
\]
The input \(K\) is recovered as a proper coordinate face.  Thus all
corner-peelable ample families still occur inside reduced forbidden
bases.  Taking \(K\) to be a cube supplies an explicit infinite family
of these bases; it does not classify all of them.
The construction also preserves every input clause change at fixed
shadow.  Proposition~\ref{prop:damp-clause-transport-minimality} uses
this to show that changing one sign bit can destroy pc-minimality while
both outputs remain nonpeelable and satisfy the reduction stopping
conditions.  Intrinsic sign admissibility for ampleness therefore does
not supply the missing minimality criterion.

\paragraph{Replicated repairs and signed-facet partitions.}
Let \(S\subseteq Q_F\) be an ample forbidden seed.  Choose a set
\(T\) of distinct missing words \(a\), such that every
\(S\cup\{a\}\) is corner-peelable and their acquired shattered
supports \(P_a\) are pairwise distinct.  On a disjoint nonempty
coordinate set \(E\), choose nonempty proper ample families
\(R_a\subsetneq Q_E\).  Form
\begin{equation}
 X=(S\times Q_E)\cup\bigcup_{a\in T}(\{a\}\times R_a).
 \label{eq:damp-introduction-common-seed}
\end{equation}
Theorem~\ref{thm:damp-common-seed-cover-classification} says that
\(X\) is forbidden exactly when each \(R_a\) is corner-peelable
and every signed facet \(\{t:t_e=i\}\) of the layer cube is
contained in at least one \(R_a\).  Fixing all layer coordinates
and intersecting the resulting sections recovers \(S\); the other
old words and their occurrences recover the tips and their families.

Call an obstruction \emph{terminal} when no corner deletion leaves
another forbidden pc-minor.  For a cornerless seed, the terminal
members of \eqref{eq:damp-introduction-common-seed} have a completely
structural description of their layer data.  Choose a surjection
\[
 \lambda:E\times\{0,1\}\longrightarrow T,
 \qquad \lambda(e,0)\ne\lambda(e,1),
\]
and set
\[
 R_a=\bigcup_{\lambda(e,i)=a}\{t\in Q_E:t_e=i\}.
\]
These conditions suffice, every terminal member in this class has this
form, and the marked sections recover \(\lambda\)
(Theorem~\ref{thm:damp-terminal-common-seed-facet-partitions}).
With two tips, the two families are cubes punctured at antipodal words.
This gives an explicit infinite family of order
\((|S|+2)2^{|E|}-2\) over every eligible two-tip seed.
The seed's forbiddenness and the resolving property of its tips remain
input conditions.

\begin{table}[tbp]
\centering
\small
\begin{tabular}{@{}p{.21\linewidth}p{.36\linewidth}p{.35\linewidth}@{}}
\toprule
Construction & What construction and recovery settle & Remaining input or coverage \\
\midrule
Alternating graph and column reductions &
Canonical reduced base; all square attachment and coordinate-block stages have intrinsic parameters and unique recovery. &
Forbidden bases of minimum degree at least three with distinct signed columns. \\
\addlinespace
Gluing at a cut vertex or separating edge &
Canonical factors and exact gluing conditions for every obstruction with the specified separator. &
Attainable sink and boundary out-map conditions of the factors; coverage outside these separator classes. \\
\addlinespace
Down-set completion &
Uniform obstruction transfer and recovery of the marked data in \eqref{eq:damp-introduction-downset}. &
The two rooted factors. \\
\addlinespace
Signed-facet assembly &
All terminal layer data over a cornerless seed with eligible tips. &
The seed and resolving tips; this construction is not exhaustive. \\
\bottomrule
\end{tabular}
\caption{Scope of the structural results.  Each row distinguishes the
parameters already constructed intrinsically from the remaining problem.}
\label{tab:damp-generative-scope}
\end{table}

\subsection{What remains to be classified}

Table~\ref{tab:damp-generative-scope} separates two outstanding tasks.
For the separator branches, the precise question is which boundary out-maps
can be realized by acyclic USOs of the factors, and how to construct and
recover supports for which the required conditions fail. This does not
replace structural generation by a search over orientations.
First, the primitive cores, rooted factors, and resolving tips need
structural admissibility conditions that imply their required peeling
properties.  Second, the construction classes need to cover all
forbidden cores.  The all-obstruction three-core reduction settles the
attachment part of the second task, but its core input is still the
original obstruction problem with a minimum-degree restriction.

Several examples delimit these tasks.  A terminal obstruction with
\(1158\) concepts and a cut vertex lies outside the replicated-repair
class; another terminal example of the same order is two-connected
and also has no such presentation
(Theorem~\ref{thm:damp-terminal-cut-vertex} and
Proposition~\ref{prop:damp-terminal-edge-amalgam}).  Thus the missing
coverage cannot be supplied merely by allowing cut vertices or by
restricting attention to two-connected terminal families.  Nor can all
obstructions be generated by corner deletions from cubes:
Corollary~\ref{cor:damp-nonpeelable-complement-obstructions} gives
an infinite family with nonpeelable cube complements, which exclude
such a descent.  Changing signs while preserving the shadow cannot always produce a
cornerless core either: a two-connected \(1162\)-concept three-core
has a shadow forcing at least three corners in every ample realization
(Corollary~\ref{cor:damp-three-core-shadow-forces-corners}).
These examples restrict proposed constructions; they do not establish
impossibility of an intrinsic classification using other finite parameters.

Nor can the one-hole construction of cornerless ample classes be
converted to a forbidden-minor construction merely by selecting its
pc-minimal outputs.  A \(580\)-concept cornerless graph three-core has
a corner in every nonempty proper pc-minor, yet its unique forbidden
minor has \(292\) concepts and one corner
(Corollary~\ref{cor:damp-cornerless-minor-minimal-not-forbidden}).
The proper-minor peeling conditions therefore retain information that
corner existence, even in all proper minors, does not supply.

The contrast with daisy cubes also concerns decomposition.  Their
peripheral coordinates expose a down-set structure from which a
generative description can be recovered.  Every ample obstruction
here is periphery-free
(Proposition~\ref{prop:damp-periphery-free}), so that starting
decomposition is unavailable.  A complete description needs a
structural replacement for the primitive cores and factors.

\subsection{Bounds, verification, and a guide to the proofs}

Let \(m_{\mathrm{amp}}\) be the least number of concepts in an ample
forbidden pc-minor.  The current bounds are
\[
 64\le m_{\mathrm{amp}}\le267
 \qquad\text{(Theorem~\ref{prop:damp-minimum-order-profile}).}
\]
The upper bound comes from an explicit family of order \(267\)
(Proposition~\ref{prop:damp-order-267-obstruction}), obtained by iterating
the contract-and-lift step that turns Hall's class into the
\(291\)-concept family of
Proposition~\ref{prop:damp-hall-order-291-obstruction}; the lower
bound combines structural reductions with finite exclusions.  Each
computationally supported result identifies its construction or proof
certificate and its checking procedure.  Uniform construction theorems
are proved symbolically; checks of fixed inputs or examples are stated
with their actual scope.  These bounds are a separate extremal result
and do not imply that the structural classification is nearly complete.

For the generative results, start with the USO and relative-extension
dictionary in Section~\ref{sec:damp-uso-boundaries}, then the basic operations and
local extension results in
Section~\ref{sec:damp-local-extension-foundations}.  Continue through
the cut-vertex and separating-edge theorems to
Section~\ref{sec:damp-canonical-attachments}, which gives the full
three-core and square-attachment proof chain before the specialized
rooted constructions.  Their dependency is explicit: the cut-vertex
theorem excludes degree-two cut vertices in an obstruction;
periphery-freeness excludes leaves; and square-deletion invariance
preserves forbidden-minor status.  These facts yield the canonical core,
and the square attachment data invert every such deletion.  The down-set construction extends the
rooted-factor examples, while the replicated-repair and signed-facet
theorems in Section~\ref{sec:damp-two-repair} describe the separate
layered construction.

For recognition, Section~\ref{sec:damp} gives the algebraic and
Yang-complex criteria.  Section~\ref{sec:damp-cores} describes
critical-core extensions, and Section~\ref{sec:damp-extensions}
develops signed minimal nonfaces and one-concept additions.
Section~\ref{sec:damp-examples} contains the fixed witnesses and the
minimum-order argument.  Appendix
\ref{app:damp-five-coordinate-relative-peeling} supplies the relative
peeling theorem used in the lower bound.  The detailed proofs and
finite certificates remain available independently of this overview;
the complete intrinsic generative classification remains open.

\input{forbidden_pc_minor_campaign/paper_sections/dismantlable_ample_classification}

\bibliographystyle{alpha}
\bibliography{main}

\end{document}
